

\label{sec:math-models}

\subsection{Synchronous Generators}
 We begin our summary of the literature by reviewing the dynamics of SGs, 
 which is a topic that has been well-established
 for decades. Despite it's lack of novelty, the topic remains highly
 relevant to the problem of IBR frequency control since
 \begin{enumerate}
  \item a substantial portion of the prevailing industry terminology
    is based around these terms, and
  \item SG dynamics interact with IBR dynamics via their network
    interconnections within the grid.
 \end{enumerate}

%  A SG is composed of three major components from a control standpoint,
% the excitation system, the 

 The relationship between power and frequency in a SG is governed by
 \begin{equation}
  J\dot{\omega} = p_m - p_e
  \label{eq:swing-equation}
 \end{equation}
 where $J$ is the system intertia, $p_m$ is the mechanical power input to
 the generator, and $p_e$ is the electrical power which is absorbed
 by connected loads.

Since a change in frequency signals there is a power imbalance within the
system, the rotor speed of an SG is typically regulated by a
feedback-mechanism called the governor. This feedback couples a change in 
frequency, $\omega$, to a (typically proportional) change in the power
delivered by the prime mover $p_m$:
\begin{equation}
  \dot{p_m} = - \frac{\omega_{set}}{R} * (\omega - \omega_{set}) + p_{m,set}
  \label{eq:governor-control}
\end{equation}

Combining equation \ref{eq:swing-equation} and \ref{eq:governor-control}
gives a second order relationship between frequency, $\omega$ and 
mechanical power delivered, $p_m$.
 
\subsection{IBR Grid Integration}
In contrast, the frequency-real power relation for IBR resources are not
coupled by mechanical inertia of the swing equation,
but rather are determined by the underlying energy source for the IBR system.
%\ref{interactivity}
defines makes the assumption that power and energy is available from a bESS,
but the same conclusions hold if the energy source is curtailed PV or wind power.
This assumption is a common and vital one: it informs the design of
the DC side of an IBR. %\ref{battery-sizing}
elaborates on the relationship between the control parameters and
the power requirements for the bESS.

%\ref{interactivity}
demonstrates that the timescale of IBR reactions are an order of magnitude
shorter than the mechanical timescale (~1 second) which SGs operate.
Thus, when considering the interaction between the two we find that 
% $p_{m,G} \approx p_{m,Go}$ is a
on this short timescale, the IBR response is first-order with respect to the
electrical power $p_{m,I}$ drawn from the circuit:
\begin{equation}
  \dot{p_{m,I}} = \frac{2 \pi (p_{e,I} - p_{m,I})}{\tau_I}
\end{equation}

It is well-documented that if the grid contains above a certain penetration
of GFLs, instability occurs. Therefore, for grids with high
IBR penetrations, a certain percentage of GFMs are necessary to maintain
stability. However, GFMs can also have their own drawbacks: they are
require more specialized hardware and control algorithms so are
therefore more expensive.

It has been shown that the placement of GFMs within a power system
topology has an important impact on overall system stability in a 
system with mixed GFL, SG, and GFM-based resources. However,
the exact interaction between GFLs and GFMs remains an open research
question. In particular, the ratio of GFLs to GFMs at which instability
arises and the optimal placement of GFMs within a given power system
remain under-studied questions.

A major difference between the physics of IBRs versus SGs is that
typically IBRs have a much lower current threshold that they
can produce based on the limits of the underlying PEs. As a result,
they have a smaller short-circuit ratio, which corresponds to 
a lower system strength. In general, a weaker system strength
implies a decrease in the voltage stiffness, which can impair
IBR performance. GFMs can provide "voltage strength" which is another
advantage over GFLs. Since this review is focused on frequency
control of IBRs, we will consider further discussion on the topic
as out of scope and will conclude with the note that voltage and
frequency control are inherently coupled and modeling these 
interactions is vital for demonstrating system stability.